% glossaries-extra v5: стиль таблицы «Условные обозначения»

\newglossarystyle{notationlong}{%
  \renewenvironment{theglossary}{%
    \small
    \setlength{\tabcolsep}{6pt}%
    \renewcommand{\arraystretch}{1.15}%
    \begin{longtable}{@{}p{0.22\linewidth}p{0.72\linewidth}@{}}%
    \toprule
    \foreignlanguage{russian}{Обозначение} & \foreignlanguage{russian}{Смысл} \\
    \midrule
    \endfirsthead
    \toprule
    \foreignlanguage{russian}{Обозначение} & \foreignlanguage{russian}{Смысл} \\
    \midrule
    \endhead
    \bottomrule
    \endfoot
  }{%
    \end{longtable}%
  }%
  \renewcommand*{\glossaryheader}{}%
  \renewcommand*{\glsgroupheading}[1]{}%
  \renewcommand{\glossentry}[2]{%
    \ifglshasdesc{##1}{%
      \glstarget{##1}{\glossentryname{##1}} &
      \glossentrydesc{##1}\glspostdescription\space ##2\tabularnewline
    }{%
      \glossentryname{##1} & \tabularnewline
    }%
  }%
  \renewcommand{\subglossentry}[3]{}%
  \renewcommand*{\glsgroupskip}{}%
}

% Запись-заголовок секции: description пустой, name = \textit{...}
\newcommand{\newnotationsection}[3]{%
  \newglossaryentry{#1}{%
    type=notation, sort={#2},
    name={#3},
    description={}%
  }%
}
